% 图 1-1 论文概览图 — 学术风格版本：平台视角而非装饰性树状图
\documentclass[border=4pt]{standalone}
\usepackage{xeCJK}
\setCJKmainfont{Songti SC}
\setCJKsansfont{Heiti SC}
\usepackage{tikz}
\usepackage{amsmath}
\usetikzlibrary{arrows.meta, positioning, calc, fit, backgrounds}
% 共享 TikZ 风格 — Aegaeon (SOSP'25) inspired，对标顶刊系统论文
% 使用：% 共享 TikZ 风格 — Aegaeon (SOSP'25) inspired，对标顶刊系统论文
% 使用：% 共享 TikZ 风格 — Aegaeon (SOSP'25) inspired，对标顶刊系统论文
% 使用：\input{_style.tex} 在 standalone preamble 内
% 调色：Okabe-Ito（与 matplotlib 图一致）

\definecolor{fuxiBlue}{HTML}{0072B2}    % 自家系统主色
\definecolor{fuxiOrange}{HTML}{E69F00}  % 第一对照
\definecolor{fuxiGreen}{HTML}{009E73}
\definecolor{fuxiPurple}{HTML}{CC79A7}
\definecolor{fuxiSky}{HTML}{56B4E9}
\definecolor{fuxiGrayBox}{HTML}{F2F2F2} % 浅灰底
\definecolor{fuxiGrayLine}{HTML}{808080}
\definecolor{fuxiGrayDark}{HTML}{555555}

\tikzset{
  % ── 通用节点 ────────────────────────────────────────
  modBox/.style = {
    draw=fuxiGrayLine, thick, fill=white,
    rounded corners=2pt, inner sep=4pt,
    minimum height=8mm, minimum width=20mm,
    align=center, font=\small
  },
  highlightBox/.style = {
    modBox, draw=fuxiBlue, fill=fuxiBlue!8,
    text=fuxiBlue, font=\small\bfseries
  },
  ghostBox/.style = {
    modBox, draw=fuxiGrayLine!50, fill=fuxiGrayBox,
    text=fuxiGrayDark, font=\small\itshape
  },
  layerBg/.style = {
    draw=fuxiGrayLine!40, dashed, fill=fuxiGrayBox!40,
    rounded corners=4pt, inner sep=8pt
  },
  externalBoundary/.style = {
    draw=fuxiGrayLine, dashed, dash pattern=on 4pt off 3pt,
    rounded corners=6pt, inner sep=10pt
  },
  % ── 箭头 ─────────────────────────────────────────────
  flow/.style = {
    -{Stealth[length=4pt, width=4pt]},
    line width=0.7pt, color=fuxiGrayDark
  },
  flowEvent/.style = {
    flow, color=fuxiBlue, line width=0.9pt
  },
  flowDup/.style = {
    flow, color=fuxiBlue!70, line width=0.7pt, dotted
  },
  bidir/.style = {
    {Stealth[length=4pt]}-{Stealth[length=4pt]},
    line width=0.7pt, color=fuxiGrayDark
  },
  % ── 编号气泡 ① ② ③ ────────────────────────────────
  numBubble/.style = {
    circle, draw=fuxiBlue, fill=white, text=fuxiBlue,
    font=\scriptsize\bfseries, inner sep=0pt,
    minimum size=4mm, line width=0.6pt
  },
}
 在 standalone preamble 内
% 调色：Okabe-Ito（与 matplotlib 图一致）

\definecolor{fuxiBlue}{HTML}{0072B2}    % 自家系统主色
\definecolor{fuxiOrange}{HTML}{E69F00}  % 第一对照
\definecolor{fuxiGreen}{HTML}{009E73}
\definecolor{fuxiPurple}{HTML}{CC79A7}
\definecolor{fuxiSky}{HTML}{56B4E9}
\definecolor{fuxiGrayBox}{HTML}{F2F2F2} % 浅灰底
\definecolor{fuxiGrayLine}{HTML}{808080}
\definecolor{fuxiGrayDark}{HTML}{555555}

\tikzset{
  % ── 通用节点 ────────────────────────────────────────
  modBox/.style = {
    draw=fuxiGrayLine, thick, fill=white,
    rounded corners=2pt, inner sep=4pt,
    minimum height=8mm, minimum width=20mm,
    align=center, font=\small
  },
  highlightBox/.style = {
    modBox, draw=fuxiBlue, fill=fuxiBlue!8,
    text=fuxiBlue, font=\small\bfseries
  },
  ghostBox/.style = {
    modBox, draw=fuxiGrayLine!50, fill=fuxiGrayBox,
    text=fuxiGrayDark, font=\small\itshape
  },
  layerBg/.style = {
    draw=fuxiGrayLine!40, dashed, fill=fuxiGrayBox!40,
    rounded corners=4pt, inner sep=8pt
  },
  externalBoundary/.style = {
    draw=fuxiGrayLine, dashed, dash pattern=on 4pt off 3pt,
    rounded corners=6pt, inner sep=10pt
  },
  % ── 箭头 ─────────────────────────────────────────────
  flow/.style = {
    -{Stealth[length=4pt, width=4pt]},
    line width=0.7pt, color=fuxiGrayDark
  },
  flowEvent/.style = {
    flow, color=fuxiBlue, line width=0.9pt
  },
  flowDup/.style = {
    flow, color=fuxiBlue!70, line width=0.7pt, dotted
  },
  bidir/.style = {
    {Stealth[length=4pt]}-{Stealth[length=4pt]},
    line width=0.7pt, color=fuxiGrayDark
  },
  % ── 编号气泡 ① ② ③ ────────────────────────────────
  numBubble/.style = {
    circle, draw=fuxiBlue, fill=white, text=fuxiBlue,
    font=\scriptsize\bfseries, inner sep=0pt,
    minimum size=4mm, line width=0.6pt
  },
}
 在 standalone preamble 内
% 调色：Okabe-Ito（与 matplotlib 图一致）

\definecolor{fuxiBlue}{HTML}{0072B2}    % 自家系统主色
\definecolor{fuxiOrange}{HTML}{E69F00}  % 第一对照
\definecolor{fuxiGreen}{HTML}{009E73}
\definecolor{fuxiPurple}{HTML}{CC79A7}
\definecolor{fuxiSky}{HTML}{56B4E9}
\definecolor{fuxiGrayBox}{HTML}{F2F2F2} % 浅灰底
\definecolor{fuxiGrayLine}{HTML}{808080}
\definecolor{fuxiGrayDark}{HTML}{555555}

\tikzset{
  % ── 通用节点 ────────────────────────────────────────
  modBox/.style = {
    draw=fuxiGrayLine, thick, fill=white,
    rounded corners=2pt, inner sep=4pt,
    minimum height=8mm, minimum width=20mm,
    align=center, font=\small
  },
  highlightBox/.style = {
    modBox, draw=fuxiBlue, fill=fuxiBlue!8,
    text=fuxiBlue, font=\small\bfseries
  },
  ghostBox/.style = {
    modBox, draw=fuxiGrayLine!50, fill=fuxiGrayBox,
    text=fuxiGrayDark, font=\small\itshape
  },
  layerBg/.style = {
    draw=fuxiGrayLine!40, dashed, fill=fuxiGrayBox!40,
    rounded corners=4pt, inner sep=8pt
  },
  externalBoundary/.style = {
    draw=fuxiGrayLine, dashed, dash pattern=on 4pt off 3pt,
    rounded corners=6pt, inner sep=10pt
  },
  % ── 箭头 ─────────────────────────────────────────────
  flow/.style = {
    -{Stealth[length=4pt, width=4pt]},
    line width=0.7pt, color=fuxiGrayDark
  },
  flowEvent/.style = {
    flow, color=fuxiBlue, line width=0.9pt
  },
  flowDup/.style = {
    flow, color=fuxiBlue!70, line width=0.7pt, dotted
  },
  bidir/.style = {
    {Stealth[length=4pt]}-{Stealth[length=4pt]},
    line width=0.7pt, color=fuxiGrayDark
  },
  % ── 编号气泡 ① ② ③ ────────────────────────────────
  numBubble/.style = {
    circle, draw=fuxiBlue, fill=white, text=fuxiBlue,
    font=\scriptsize\bfseries, inner sep=0pt,
    minimum size=4mm, line width=0.6pt
  },
}


\begin{document}
\begin{tikzpicture}[font=\sffamily, node distance=8mm and 13mm]

% ── 主执行闭环 ─────────────────────────────────────
\node[modBox, minimum width=26mm, minimum height=12mm] (user) {
  用户入口\\[-1pt]\scriptsize IM/PWA · REPL
};

\node[highlightBox, right=16mm of user, minimum width=30mm, minimum height=13mm] (xuannv) {
  玄女编排层\\[-1pt]\scriptsize intent · task · state
};

\node[modBox, right=16mm of xuannv, minimum width=30mm, minimum height=13mm] (a2a) {
  A2A 协议适配\\[-1pt]\scriptsize JSON-RPC · SSE
};

\node[modBox, right=16mm of a2a, minimum width=30mm, minimum height=13mm] (workers) {
  门客执行层\\[-1pt]\scriptsize Claude Code · Codex
};

% ── 事件与审计平面 ─────────────────────────────────
\node[highlightBox, below=16mm of a2a, minimum width=42mm, minimum height=13mm] (bus) {
  EventBus 事实流\\[-1pt]\scriptsize live broadcast + WAL append
};

\node[ghostBox, below left=10mm and 8mm of bus, minimum width=30mm, minimum height=11mm] (wal) {
  SQLite WAL\\[-1pt]\scriptsize replay · audit
};

\node[modBox, below right=10mm and 8mm of bus, minimum width=30mm, minimum height=11mm] (obs) {
  多端观测\\[-1pt]\scriptsize Firehose · SSE · IM
};

% ── 执行隔离 ───────────────────────────────────────
\node[ghostBox, below=10mm of workers, minimum width=30mm, minimum height=11mm] (sandbox) {
  三层沙箱\\[-1pt]\scriptsize L1 · L2 worktree · L3
};

% ── 主控制流 ───────────────────────────────────────
\draw[flow] (user.east) -- node[above=1pt, font=\scriptsize] {任务/人工回复} (xuannv.west);
\draw[flow] (xuannv.east) -- node[above=1pt, font=\scriptsize] {dispatch} (a2a.west);
\draw[flow] (a2a.east) -- node[above=1pt, font=\scriptsize] {send/stream} (workers.west);
\draw[flow] (workers.south) -- node[right=1pt, font=\scriptsize] {隔离执行} (sandbox.north);

% ── 事件流：汇入 EventBus，再分发/持久化 ──────────────
% 门客 → bus：起点下移 4mm 避开灰色 send/stream 箭头，左拐再下到 bus 右沿
\draw[flowEvent] ($(workers.west)+(0,-4mm)$) -- ++(-6mm,0) |- ($(bus.east)+(0,3mm)$);
\draw[flowEvent] (xuannv.south) |- ($(bus.west)+(0,3mm)$);

\draw[flowEvent] (bus.south west) -- (wal.north east);
\draw[flowDup] (bus.south east) -- (obs.north west);

% ── input-required 回环 ────────────────────────────
\draw[flow, color=fuxiPurple] (bus.west) -| ($(xuannv.south)+(-8mm,0)$)
  node[pos=0.18, below=1pt, font=\scriptsize, color=fuxiPurple] {input-required};
% 玄女 → 用户：直接一条平行短线（不再折角），起点在 xuannv 左沿下 4mm
\draw[flow, color=fuxiPurple] ($(xuannv.west)+(0,-4mm)$) -- ($(user.east)+(0,-4mm)$);

% ── 背景分区 ───────────────────────────────────────
\begin{scope}[on background layer]
  \node[layerBg, fit=(user) (xuannv) (a2a) (workers), inner sep=5pt,
    label={[font=\scriptsize\itshape, text=fuxiGrayDark]above:{控制平面：从用户意图到门客执行}}] {};
  \node[layerBg, fit=(bus) (wal) (obs), inner sep=5pt,
    label={[font=\scriptsize\itshape, text=fuxiGrayDark]below:{事实平面：可观测、可回放、可审计的事件序列}}] {};
\end{scope}

% ── 图例 ───────────────────────────────────────────
\node[anchor=north east, font=\scriptsize, align=left, draw=fuxiGrayLine!40, fill=white,
      inner sep=4pt, line width=0.3pt] at ($(workers.north east)+(0,11mm)$) {
  \textcolor{fuxiGrayDark}{$\rightarrow$ 控制流 \quad}
  \textcolor{fuxiBlue}{$\rightarrow$ 事件流 \quad}
  \textcolor{fuxiPurple}{$\rightarrow$ 人工介入回环}
};

\end{tikzpicture}
\end{document}
